\section{Related Work}
\label{sec:related}

\subsection{Efficient Architectures and Low-Resolution Vision}
The evolution of computer vision architectures has largely been driven by the quest for efficiency and scalability. Tan and Le introduced \textbf{EfficientNetV2} \cite{tan2021efficientnetv2smallermodelsfaster}, highlighting the importance of training speed and parameter efficiency through the use of Fused-MBConv layers. While highly effective on ImageNet, standard CNNs often struggle with information loss when applied to small inputs. Addressing this, Hassani et al. proposed \textbf{Compact Convolutional Transformers (CCT)} \cite{hassani2022escapingbigdataparadigm}, demonstrating that transformers can escape the ``big data paradigm'' by using convolutional tokenizers to handle small datasets effectively. Similarly, modern pure ConvNets like \textbf{ConvNeXt V2} \cite{woo2023convnextv2codesigningscaling} utilize fully convolutional masked autoencoders to improve representation learning. More recently, \textbf{Vision Mamba (Vim)} \cite{zhu2024visionmambaefficientvisual} challenged the transformer dominance by utilizing bidirectional State Space Models (SSMs) for linear complexity scaling, offering a computationally efficient alternative for visual representation. Our work evaluates these state-of-the-art architectures, adapting their downsampling strategies to the constraints of $64 \times 64$ facial expression recognition (FER).

\subsection{Transfer Learning in Facial Analysis}
While generic object recognition models provide a strong baseline, FER benefits significantly from domain-specific pre-training. Transfer learning from massive face recognition datasets (e.g., CelebA, MS-Celeb-1M) is a standard practice to overcome the data scarcity in emotion recognition. Research indicates that models pre-trained on millions of faces learn robust, identity-invariant features that are highly transferable. Wang et al. demonstrated this in their \textbf{QCS (Quadruplet Cross Similarity)} network \cite{wang2025qcsfeaturerefiningquadruplet}, utilizing an InsightFace ResNet-50 (IR50) backbone pre-trained on MS-Celeb-1M to extract fine-grained facial features. Inspired by their finding that pre-trained FR backbones provide superior initialization, we adopt the IR50 weights from their repository as a starting point. However, we address a critical limitation: the massive classification heads designed for thousands of identities lead to rapid overfitting on smaller emotion datasets. We propose replacing these heavy heads with lightweight, attention-based mechanisms to preserve spatial information without parameter bloat.

\subsection{Lightweight Models and Transfer Learning Strategies for Efficient FER}

Recent advances in facial expression recognition (FER) have largely been driven by increasingly deep and parameter-heavy convolutional neural networks. However, the CERN \cite{GERA20229} architecture challenges this trend by demonstrating that state-of-the-art (SOTA) performance can be achieved with only 1.45 million parameters. Despite its compact design, CERN outperformed several larger and more computationally demanding models across benchmark datasets. On in-the-wild datasets RAFDB and FerPlus CERN reports 87.09% and 88.17% accuracies.


\subsection{Optimization and Generalization Techniques}
Training deep networks on heterogeneous, imbalanced datasets requires specialized optimization strategies to ensure convergence and generalization. 
\begin{itemize}
    \item \textbf{Warm-up and Learning Rate Schedules:} As demonstrated by Goyal et al. \cite{goyal2018accuratelargeminibatchsgd}, utilizing a learning rate warm-up period is critical for stabilizing training, particularly when using large minibatches. This technique prevents the model from diverging due to large gradient updates during the initial epochs when weights are random.
    \item \textbf{Sharpness-Aware Minimization (SAM):} Standard optimizers like AdamW focus solely on minimizing the training loss value. Foret et al. introduced \textbf{SAM} \cite{foret2021sharpnessawareminimizationefficientlyimproving}, which simultaneously minimizes loss value and loss sharpness. By seeking parameters that lie in ``flat minima,'' SAM significantly improves generalization performance on out-of-distribution data, a crucial property for our multi-source dataset that includes both lab-controlled and in-the-wild images.
\end{itemize}

\subsection{Choosing a FACS-Aligned Benchmark Dataset for Robust Generalization Testing}

CK+ is a well-established dataset focusing on prototypical and micro-expressions, with annotations grounded in the Facial Action Coding System (FACS) developed by Ekman and Friesen~\cite{Ekman1978FACS}. FACS defines facial movements in terms of Action Units (AUs), which correspond to specific underlying muscle activations. The selection of key facial landmarks and their configuration in CK+ follows the FACS framework~\cite{Lucey2010CK}, making it particularly suitable for evaluating models designed to capture fine-grained facial expression dynamics.

